Pour chaque acquisition, il existe un fichier contenant toutes
ses méta-données.

Il est utilisable par les logiciels de visualisation et traitement,
ou simplement lisible par l'utilisateur. Il facilite le travail
de dépouillement des données.

Le format de ce fichier est décrit dans ce chapitre.

\subsection{Nom du fichier}

Pour une acquisition donnée le fichier ``inf'' s'appelle :
\texttt{\$\{Acq\_Id\}.inf}

 Acq\_Id est l'identificateur d'acquisition AIDA. Par exemple,
\path{fa110015} pour la quinzième acquisition de l'année 2011
sur le site du Fauga.
    Il n'y a aucune correspondance entre le numéro et la date de mesure.


\subsection{Stockage}

Tous les fichiers sont rangés dans le répertoire : 

\texttt{\$AIDA\_META\//INF}


\subsection{Format}

C'est un fichier de type texte, constitué d'une série de paragraphes,
eux-mêmes constitués d'une série d'égalités "champ = valeur". 

Le début d'un paragraphe est défini par la présence de son
nom en majuscules et entre crochets suivi d'un retour à la ligne.

A l'intérieur d'un paragraphe, la fin de ligne est marquée
par un point virgule:~\texttt{;}.

A l'intérieur de la valeur les informations sont séparées par un caractère:
\texttt{/}. Ces informations ne doivent pas contenir ce caractère \texttt{/}.

Les lignes vides, les espaces, et les tabulations sont permis.

Les lignes de commentaires doivent commencer par un caractère
dièse: \texttt{\#};

Les champs ne doivent pas contenir d'accents.

Les valeurs peuvent contenir des accents, à condition d'utiliser
le codage iso8859-1 (ou latin1).

Le fichier doit se terminer par \texttt{{[}FIN{]}}.

La liste des paragraphes ainsi que leur caractère "obligatoire"
ou "optionnel" est présentée dans le tableau ci-dessous.

\begin{center}

\begin{tabular}{|l|l|}
\hline 
Nom du Paragraphe& Caractère\\
\hline 
\hline 
\textbf{FORMAT}& \textbf{obligatoire}\\
\hline 
\textbf{IDENTIFICATION}& \textbf{obligatoire}\\
\hline 
\textbf{PARAMETRES\_ACQUIS}&\textbf{obligatoire}\\
\hline 
\textbf{ACQUISITION}& \textbf{obligatoire}\\
\hline 
HORAIRES\_UTC & optionnel\\
\hline 
INS& optionnel\\
\hline 
SEQUENCES& optionnel\\
\hline 
PERSONNEL& optionnel\\
\hline 
POINTS\_CARACTERISTIQUES& optionnel\\
\hline 
METEO\_DEBUT& optionnel\\
\hline  
METEO\_FIN& optionnel\\
\hline 
ETAT\_CAPTEURS& optionnel\\
\hline 
\end{tabular}
\end{center}


\subsubsection{{[}FORMAT{]}}

\label{parag-FORMAT}
Ce paragraphe décrit les formats utilisés dans les valeurs.

\ttfamily
\saute
\noindent\begin{tab2col}
\hline 
\textnormal{Champ}& \textnormal{Valeur}\\
\hline 
\hline 
format\_coord &
C (D)DDdMM.X(X) (C=lettre (N, E, S ou W), (D)DD=degrés 2 ou 3 chiffres, MM=minutes,
X(X)=fraction décimale de minutes;\\
\hline 
format\_heure&
HHMMSS (HH=heures, MM=minutes, SS=secondes);\\
\hline 
format\_date&
JJ/MM/AAAA (JJ=jour, MM=mois, AAAA=année, le tout en chiffres);\\
\hline 
format\_duree&
HH.D (HH=heure, D=dixièmes d'heure);\\
\hline 
\end{tab2col}
\saute

\normalfont

\subsubsection{{[}IDENTIFICATION{]}}
\label{parag-IDENTIFICATION}


\saute
\noindent\begin{tab4col}
\hline 
Champ& Caractère& Format valeur& Description \\
\hline 
\hline 
\texttt{plate-forme}& \textbf{obligatoire} & \textit{texte;}&
nom de la plate-forme en toutes lettres \\
\hline 
\texttt{acq\_id}& \textbf{obligatoire} & \textit{texte;} & identificateur d'acquisition AIDA\\
\hline 
\texttt{date}& \textbf{obligatoire} & \textit{JJ/MM/AAAA;} &
JJ= jour, MM= mois, AAAA= année (format décrit dans le paragraphe FORMAT)\\
\hline 
\texttt{campagne}& \textbf{obligatoire} & \textit{texte;}& nom de la campagne\\
\hline 
\texttt{theme}& optionnel& \textit{texte;}& thème scientifique de l'acquisition\\
\hline 
\texttt{lieu\_depart}& optionnel& \textit{texte;}&
lieu de départ de la plate-forme de mesure\\
\hline 
\texttt{lieu\_arrivee}& optionnel& \textit{texte;}&
lieu d'arrivée de la plate-forme de mesure\\
\hline 
\end{tab4col}
\saute

 Le nom de la plate-forme doit être connu d'AIDA par le fichier
\texttt{plate-forme}

Le champ \texttt{campagne} doit contenir une seule valeur.
Les acquisitions faites
hors de la période opérationnelle de la campagne ont comme nom de
campagne ESSAI, CALIB, ou TRANSIT.

Le champ \texttt{theme} doit expliciter plus librement le but de 
l'acquisition (ex : essais réfractomètre, vol type P1-flux, 
transit Brétigny-Dole, ....).

\subsubsection{{[}HORAIRES\_UTC{]}}

Ce paragraphe est utilisé dans le cas où la plate-forme de mesure est "mobile"
(avion, ballon, drone à ailes...)

\newcommand{\formathoraire}{\textit{hhmmss} ou\newline \textit{AAAA}-\textit{MM}-\textit{JJ}T\textit{hh}:\textit{mm}:\textit{ss}}
\saute
\noindent\begin{tab3col}
\hline 
Champ& Format valeur& Description\\
\hline 
\hline 
\texttt{depart\_parking}& \formathoraire & heure UTC de départ du parking\\
\hline 
\texttt{decollage}& \formathoraire & heure UTC de décollage\\
\hline 
\texttt{atterrissage}& \formathoraire & heure UTC d'atterrissage\\
\hline 
\texttt{retour\_parking}& \formathoraire & heure UTC de retour au parking\\
\hline 
\texttt{duree\_reglementaire}& \textit{hh.d}& durée réglementaire du vol\\
\hline 
\end{tab3col}
\saute


L'existence de deux formats impose aux programmes qui utilisent
ces données de savoir décoder les deux. La distinction se fait sur la
longueur du champ. La forme courte n'indique qu'une heure, et est donc ambigüe
dans les cas où l'acquisition s'étend sur plusieurs jours. Les différents cas
possibles sont :
\begin{itemize}
 \item utilisation d'un champ \textit{hh} >=24. Dans ce cas, l
 \item utilisation d'un champ \textit{hh} <24, mais avec une fin 
antérieure au début (par exemple \texttt{decollage = 223344} et 
\texttt{atterrissage = 001122})
\end{itemize}


Pour l'écriture de 
La forme longue du champ est recommandée, car elle n'est pas ambigüe,
et suit un format reconnu (ISO8601) par de nombreux langages et 
bibliothèques informatiques.

On reconnait que, lorsque les données sont saisies manuellement, 
la forme courte est plus concise et rapide. Il est cependant recommandé 
d'appliquer rapidement une normalisation après


\subsubsection{{[}ACQUISITION{]}}

\saute
\noindent\begin{tab4col}
\hline 
Champ& Caractère & Format valeur& Description \\
\hline 
\hline 
\texttt{systeme\_acquisition}& \textbf{obligatoire} & \textit{texte;}& nom du système d'acquisition\\
\hline 
\texttt{debut\_acquisition}& \textbf{obligatoire} & \textit{JJ/MM/AAAA/hhmmss;}& 
date et heure de début d'acquisition \\
\hline 
\texttt{fin\_acquisition}& \textbf{obligatoire} & \textit{JJ/MM/AAAA/hhmmss;}&
date et heure de fin d'acquisition \\
\hline 
\texttt{media}& \textbf{obligatoire} & \textit{texte;}& 
support de stockage du système d'acquisition\\
\hline 
\texttt{autre\_acquisition}& optionnel & \textit{texte/texte/...;}&
nom d'autres acquisitions sur la m\^eme plate-forme\\
\hline 
\texttt{commentaire}& optionnel & \textit{texte;}& commentaire\\
\hline 
\end{tab4col}
\saute

\noindent\texttt{media} peut prendre les valeurs floppy, dat, exabyte, hd,
don

\texttt{systeme\_acquisition} peut prendre les valeurs :
\begin{itemize}
\item RESUM, CL, PP, PPI pour l'ARAT
\item MC, MC95  pour le Merlin IV
\item EDISOL, ... pour les stations sols 
\item 4M-BATEAU, .... pour les acquisitions 4M  sur bateau
\item STAR, ... pour les radio-sondages
\end{itemize}

\subsubsection{{[}INS{]}}

Ce paragraphe donne les principales informations relatives à
la centrale à inertie de la plate-forme.

\saute
\noindent\begin{tab3col}
\hline 
Champ& Format valeur& Description\\
\hline 
\hline 
\texttt{ramp\_depart}& \textit{Nxxdxx.x/Wxxxdxx.x/xxxxft;}&
coordonnées du point de départ\\
\hline 
\texttt{type\_aln}&
\textit{xxx}~\texttt{\textit{s}}\textit{;} ou \textit{xxx}~\texttt{\textit{mn}}\textit{;}&
durée de l'alignement en secondes ou minutes\\
\hline 
\texttt{ramp\_arrivee}&
\textit{Nxxdxx.x/Wxxxdxx.x/xxxxft;}& coordonnées du point d'arrivée\\
\hline 
\texttt{pp\_arrivee}& \textit{Nxxdxx.x/Wxxxdxx.x/xxxxft;}&
coordonnées réelles à l'arrivée\\
\hline 
\end{tab3col}
\saute

\subsubsection{{[}POINTS\_CARACTERISTIQUES{]}}

Ce paragraphe décrit les différents points 
géographiques intéressant l'acquisition.

\saute
\noindent\begin{tab3col}
\hline 
Champ& Format valeur& Description\\
\hline 
\texttt{wpt}& \textit{nom/Nxxdxx.x/Wxxxdxx.x;}&
abréviation et coordonnées du point \\
\hline 
\texttt{radar}& \textit{nom/Nxxdxx.x/Wxxxdxx.x;}&
abréviation et coordonnées du point \\
\hline 
\texttt{balise}& \textit{nom/Nxxdxx.x/Wxxxdxx.x;}&
abréviation et coordonnées du point \\
\hline 
\texttt{bouee}& \textit{nom/Nxxdxx.x/Wxxxdxx.x;}&
abréviation et coordonnées du point \\
\hline 
\texttt{station}& \textit{nom/Nxxdxx.x/Wxxxdxx.x;}&
abréviation et coordonnées du point \\
\hline 
\texttt{ville}& \textit{nom/Nxxdxx.x/Wxxxdxx.x;}&
abréviation et coordonnées du point \\
\hline 
\end{tab3col}
\saute

\subsubsection{{[}PERSONNEL{]}}

\saute
\noindent\begin{tab3col}
\hline 
Champ&
Format valeur&
Description\\
\hline 
\hline 
\texttt{pilote}& \textit{Initiale du prénom. Nom;}&
Nom du pilote de la plate-forme (avion)\\
\hline 
\texttt{copilote}& idem&
Nom du copilote de la plate-forme\\
\hline 
\texttt{operateur}& idem&
Nom de l'opérateur du système d'acquisition\\
\hline 
\texttt{scientifique}& idem&
Nom du scientifique de la mission\\
\hline 
\texttt{directeur\_de\_mission}& idem&
Nom du directeur de la mission\\
\hline 
\texttt{invite}& idem&
Nom de l'invité \\
\hline 
\texttt{commandant\_de\_bord}& idem&
Nom du commandant de bord (bateau)\\
\hline 
\end{tab3col}
\saute

Les champs peuvent être répétés si plusieurs personnes
correspondent à la même désignation (ex: plusieurs 
scientifiques).


\subsubsection{{[}METEO\_DEBUT{]}}

Ce paragraphe décrit les conditions météorologiques 
en début d'acquisition.

\saute
\noindent\begin{tab3col}
\hline 
Champ& Format valeur& Description\\
\hline 
\hline 
\texttt{QNH}& \textit{xxx hPa;} ou \textit{xxx inHg;} 
ou  \textit{xxx mmHg;}& Valeur du QNH\\
\hline 
\texttt{QFE}& \textit{xxx hPa;} ou \textit{xxx inHg;} 
ou  \textit{xxx mmHg;} & Valeur du QFE\\
\hline 
\texttt{vent}& \textit{ddd/ff kt;} ou  \textit{ddd/ff m.s-1;}
& Direction en degrés et force du vent réel\\
\hline 
\texttt{temperature}& \textit{xxx C;} ou \textit{xxx F;} & Température en degrés C ou F\\
\hline 
\texttt{point\_de\_rosee}& \textit{xxx C;} ou \textit{xxx F;} & Température en degrés C ou F\\
\hline 
\texttt{visibilite}& \textit{xxxx m;} ou \textit{xx km;}& Visibilité\\
\hline 
\texttt{humidite}& \textit{xx \%;} & Humidité relative\\
\hline
\texttt{nuage}& \textit{texte;} & Description des nuages\\
\hline
\texttt{temps\_present}& \textit{texte;} & Description du temps présent\\
\hline
\texttt{origine}& \textit{texte;} & Origine de l'information météorologique\\
\hline
\texttt{commentaire}& \textit{texte;} & Commentaire de l'opérateur\\
\hline
\end{tab3col}
\saute

Le champ commentaire sert à recueillir les remarques de
l'opérateur sur l'information météorologique. Celles-ci
peuvent par exemple, décrire l'évolution rapide d'une situation,
ou apporter des informations non fournies dans les autres champs
du paragraphe.

\subsubsection{{[}METEO\_FIN{]}}

Ce paragraphe décrit les conditions météorologiques en fin
d'acquisition.

Le format est identique à \hbox{[METEO\_DEBUT]}.


\subsubsection{{[}SEQUENCES{]}}

Ce paragraphe décrit les différentes séquences spatiales et temporelles
au cours de l'acquisition.

\saute
\noindent\begin{tab2col}
\hline 
Champ& Format valeur\\
\hline 
\texttt{palier}&
\textit{nom/hhmmss\( _{deb} \)/hmmss\( _{fin} \)/Wpt\( _{deb} \)/Wpt\( _{fin} \)/Cap/Alt./comment;}\\
\hline 
\texttt{sondage}&
\textit{nom/hhmmss\( _{deb} \)/hmmss\( _{fin} \)/Wpt\( _{deb} \)/Wpt\( _{fin} \)/Alt\( _{deb} \)/Alt\( _{fin} \)/comment;}\\
\hline 
\texttt{cercle}&
\textit{nom/hhmmss\( _{deb} \)/hmmss\( _{fin} \)/Wpt/Alt/comment;}\\
\hline 
\texttt{zoom}&
\textit{nom/hhmmss\( _{deb} \)/hmmss\( _{fin} \)/Wpt\( _{deb} \)/Wpt\( _{fin} \)/Cap/Alt\( _{deb} \)/Alt\( _{fin} \)/comment;}\\
\hline 
\texttt{virage}&
\textit{nom/hhmmss\( _{deb} \)/hmmss\( _{fin} \)/Cap\( _{av} \)/Cap\( _{ap} \)/Angle parcouru/Alt/comment;}\\
\hline 
\texttt{station}&
\textit{nom/hhmmss\( _{deb} \)/hmmss\( _{fin} \)/Wpt/Cap/Alt/comment;}\\
\hline 
\texttt{sequence}&
\textit{nom/hhmmss\( _{deb} \)/hmmss\( _{fin} \)/comment;}\\
\hline 
\end{tab2col}
\saute

 \textit{nom} : abréviation de la séquence

 \textit{hhmmss\( _{deb} \)} : \textit{}heure de début de la
séquence sous la forme hhmmss

 \textit{hmmss\( _{fin} \)} : \textit{}heure de début de la
séquence sous la forme hhmmss

 \textit{Wpt} : abréviation du point géographique où se déroule
la séquence

 \textit{Wpt\( _{deb} \)} : abréviation du point de début de la
séquence

 \textit{Wpt\( _{fin} \)} : abréviation du point de fin de la
séquence

 \textit{Cap} : cap de la plate-forme durant la séquence en
degrés

 \textit{Cap$_{av}$} : cap de la plate-forme avant la séquence en
degrés 

 \textit{Cap$_{ap}$} : cap de la plate-forme après la séquence en
degrés

 \textit{Alt.} : altitude de la plate-forme durant la séquence
suivie de l'unité ft ou m, ou notée en niveau de vol FLxxx

 \textit{Alt\( _{deb} \)} : altitude de la plate-forme au début
de la séquence suivie de l'unité ft ou m, ou notée en niveau de vol FLxxx

 \textit{Alt\( _{fin} \)} : altitude de la plate-forme au début
de la séquence suivie de l'unité ft ou m,
ou notée en niveau de vol FLxxx

 \textit{Angle parcouru} : angle parcouru par la plate-forme au cours de la
séquence. Il est noté positif dans le sens des aiguilles d'une montre.

 \textit{comment} : commentaire texte


\subsubsection{{[}ETAT\_CAPTEURS{]}}

\saute
\noindent\begin{tab3col}
\hline 
Champ& Format valeur& Description\\
\hline 
\hline 
\texttt{avant\_acq}& \textit{texte;} &
commentaire sur l'état des capteurs avant l'acquisition\\
\hline 
\texttt{pendant\_acq}& \textit{texte;} &
commentaire sur l'état des capteurs pendant l'acquisition\\
\hline 
\texttt{apres\_acq}& \textit{texte;} &
commentaire sur l'état des capteurs après l'acquisition\\
\hline 
\end{tab3col}
\saute


\subsubsection{{[}PARAMETRES\_ACQUIS{]}}

\saute
\noindent\begin{tab3col}
\hline 
Champ& Format valeur& Description valeur\\
\hline 
\hline 
\texttt{abréviation}& \textit{texte;}& nom détaillé du paramètre\\
\hline 
\end{tab3col}
\saute

Ce paragraphe contient une ligne "\texttt{champ = valeur}" par paramètre
acquis. 
Le nom du champ est l'abréviation du paramètre.\\ 
Ce paragraphe peut être très long, car pour certaines plates-formes
le nombre de paramètres est important. Par conséquent pour améliorer la lisibilité
du fichier, ce paragraphe doit être \textbf{placé en fin de fichier}.


